 % !TeX root = book.tex
\chapter{Introduction}
\label{sec:intro}

\vspace{-4.5cm}

\begin{figure*}[h]
    \centering
	\includegraphics[width=6cm]{spiders-web-2.png}
\end{figure*}

\vspace{-0.5cm}

\noindent 
This book articulates a \textit{vision} -- an ambitious and forward-looking proposal -- for a reconceptualization of Prolog in the context of the modern Web. At the heart of this vision lies \textit{Web Prolog}, a carefully crafted profile of the Prolog language, reimagined as a logic-based programming language for the Web. More than just a variant of Prolog adapted for networked environments, Web Prolog aspires to function as a foundational language for a new generation of \textit{web-native logic programming} and \textit{web-based agent systems}.

We propose to think of Web Prolog not only as a tool for building programs that interact over the Web, but as a language for defining logic-driven agents -- \textit{Prolog agents} -- capable of reasoning, communicating, and coordinating their behavior across distributed contexts. (The term ``agent'' is used very liberally throughout this book and does not imply agency in its fullest sense.) These agents inhabit a conceptual space we call the \textit{Prolog Web}: an architectural extension of the traditional Web, designed to accommodate structured logical interactions, intelligent agent communication, and declarative knowledge sharing.

Taken together, these three components form the \textit{Prolog Trinity ecosystem}. This ecosystem is envisioned as a niche within the broader World Wide Web, offering a unique paradigm for symbolic AI, declarative communication, and logical coordination on the Web. It represents not merely a technical enhancement, but a philosophical stance: a reassertion of the relevance of symbolic reasoning and declarative knowledge in an era increasingly dominated by opaque statistical models.

This opening chapter serves as a prologue to that larger vision. It prepares the ground for the conceptual foundations, technical proposals, and speculative possibilities explored in the chapters ahead, and it introduces the \textit{what}, \textit{why} and \textit{how} of a project that aims to bring it to life.

\vspace{-0.5cm}
\section{The Prolog Trinity ecosystem}

\begin{quotation}
Employed as a language for communicating with computers, logic is higher-level and more human-oriented than other formalisms specifically developed for computers.
\vspace{0cm}
\flushright \textit{Robert Kowalski}
\end{quotation}

\noindent  
As suggested by the simple diagram in Figure~\ref{fig:the-big-picture-4}, we think of Web Prolog, Prolog agents and the Prolog Web as three different kinds of entities, playing three distinct roles in the comprehensive Prolog-based ecosystem that we think of as the Trinity. 

\begin{figure}[h]
    \centering
	\includegraphics[width=5.5cm]{trinity}
    \caption{The Prolog Trinity ecosystem: Web Prolog, Prolog agents, and the Prolog Web, and their mutual dependencies.}
    \label{fig:the-big-picture-4}
\end{figure}
\vspace{-0.2cm}
\noindent The diagram illustrates the main idea on the highest level of abstraction we shall care about. The arrows are meant to indicate the many interdependencies and symbiotic relationships that exist between these entities, and to remind ourselves that we are looking at one and the same ecosystem from three different angles. 

Put simply, on this level of abstraction the Prolog Web is a network of Prolog agents being programmable in Web Prolog and talking Web Prolog to each other. Expressed more elaborately, the Prolog Web is populated, and indeed also propped up, by Prolog agents, and provides, in the form of distributed Web Prolog databases, some of the knowledge that these agents need to do their thing. Prolog agents in the Prolog Web environment interact with other Prolog agents, using Web Prolog as their language of communication. They are also able to talk with (other kinds of) agents outside this environment such as humans, maybe through a conventional graphical user interface (a GUI), or perhaps using spoken natural language in a voice user interface (a VUI), and maybe even in a dialog with an \textit{embodied conversational agent} (an ECA).

It is important to understand that the Prolog Trinity ecosystem is not an island. The Prolog Web does not exist in complete isolation, cut off from the rest of the Web. Rather, it is meant to serve as an \textit{extension} of the traditional Web, along with other extensions such as the Semantic Web and the Web of Things. Just like these, the Prolog Trinity can be regarded as a \textit{sub-ecosystem} of the Web ecosystem in its entirety. As we shall see, it can also be thought of a sub-ecosystem of the larger ecosystem surrounding Prolog -- the already existing Prolog systems and the communities of logic programming researchers,  system implementers, and other users of Prolog.

\vspace{-0.5cm}
\section{Walking in the footsteps of inventors}
\vspace{-0.2cm}
Before we continue, let us first pay tribute to five people who invented technologies that form the basis for the work presented in the book and laid the groundwork for the Prolog Trinity ecosystem. We have put them side by side in Figure~\ref{fig:inventors}. 

\begin{figure}[h]
    \centering
	\includegraphics[width=\textwidth]{inventors}
    \vspace{-0.3cm}\caption{Five inventors whose ideas underpin the Trinity: Colmerauer, Kowalski, Armstrong, Berners-Lee, and Harel (photo montage).}
    \label{fig:inventors}
\end{figure}

\vspace{-0.3cm}
\noindent On the left we find Alain Colmerauer. He invented Prolog in 1972 and built the first implementation. Sadly, he passed away in 2017. To the right of him is Robert Kowalski -- who came up with a lot of the theory behind Prolog which inspired Colmerauer. He is still around, and is still doing a lot of logic programming research and development. Third from the left, we find Joe Armstrong, who invented another programming language called Erlang. He liked Prolog, and to some extent, as we shall see, Erlang was inspired by Prolog. Regretfully, Joe Armstrong died in 2019. To the right of him, we find Tim Berners-Lee. He invented the Web, the first web browser, and the protocols and algorithms allowing the Web to scale. He is still very much alive and active trying to build a better Web. Finally, on the right, we find David Harel, the inventor of Statecharts, a rather succesful visual state-machine programming language. Why his work has a role to play in the Prolog Trinity ecosystem will be revealed later in the book. 

These five inventors are (or were, for those that are no longer with us) scientists and engineers who have made significant contributions to the field of computer science. We do not mean to suggest that they were alone in inventing these technologies.\footnote{For example, Robert Virding and Mike Williams were heavily involved in the design and development of Erlang.} The work presented in this book is indebted also to the many, many scholars who evolved the original ideas into the shape they have today. This book tries to build on their achievements, combining chunks of their work into a what we hope will emerge as a coherent whole, and preferably a whole greater than the simple sum of its parts.


\section{The Prolog Trinity ecosystem concretized}\label{sec:trinity-concretized}

A deeper understanding of what the Prolog Trinity is all about can hopefully be gained by considering the tiny partition of the Prolog Web suggested in Figure~\ref{fig:prolog-web-intro}. 

\begin{figure}[h]
    \centering
	\includegraphics[width=\textwidth]{prolog-web-intro-10.png}
    \caption{End-user clients interacting with Prolog agents within a tiny partition of the Prolog Web.}
    \label{fig:prolog-web-intro}
\end{figure}

\noindent On this fairly low level of abstraction, the Prolog Trinity ecosystem can be understood as a concrete network of concrete \textit{Prolog nodes}. In Figure~\ref{fig:prolog-web-intro}, nodes are shown as boxes labeled by a \textit{domain name} at the top. The arrows symbolize connections between different nodes, and between end-user clients and nodes. As suggested by the diagram, each node is divided into two ``containers,'' one container dedicated to the management of actor processes, the other to           the storage of Web Prolog programs and data. 

\subsection{Prolog nodes and Prolog actors}

Inspired by Erlang's notion of an actor, the \textit{Prolog actor} is the fundamental unit of computation -- and the simplest type of agent -- in the Prolog Trinity ecosystem. The smiley emojis in Figure~\ref{fig:prolog-web-intro} represent Prolog actors of various kinds that live and work at nodes. When a node is computing something for a client, it is often one or more of these actors that are responsible, even if this is not evident from the outside. Hiding an actor behind an emoji makes it look simpler than it actually is, and in fact, an actor has an interesting internal structure. The diagram in Figure~\ref{fig:actor-0} shows three actors, one of which has just been spawned and one that has been opened up so that we can look inside (where, at this level of abstraction, they all look the same).

\begin{figure}[h]
    \centering
	\includegraphics[width=5.6cm]{an-actor}
    \caption{The anatomy of a Prolog actor.}
    \label{fig:actor-0}
\end{figure}

\noindent A Prolog actor is a \textit{process}, capable of communicating with the world around it through messaging. An actor has a unique address -- a \textit{process identifier}, or \textit{pid} for short. If we have the pid of an actor, we can message it, and since pids can themselves be part of a message, it allows other actor processes to communicate back. Furthermore, an actor has a \textit{mailbox} that stores incoming messages. There is a \textit{receive} operation that allows the actor to select and process messages from the mailbox, and a \textit{send} operation that can be used to send messages to any other actor. There is a \textit{spawn} operation that allows an actor to create other actors -- child actors as it were. Finally, every actor is equipped with a private Prolog clause database, a dedicated space where programs and data exclusive to that actor are stored. 

Authorized end-user clients can connect and interact with nodes. These clients may range from web browsers running simple JavaScript-based Prolog terminals or intelligent web applications, to more sophisticated devices like VR headsets or social robots. (Note that not all clients are end-user clients because nodes can be clients too.) Human users interacting with end-user clients typically trigger activity in the Prolog Web, where agents as well as connections between agents are more or less frequently being created, interacted with, and then destroyed. Some agents and some connections live for a long time, while others may only last for fractions of a second. This is obviously a form of dynamics that is not easy to capture in a static diagram such as Figure~\ref{fig:prolog-web-intro}.

A node can be seen as a special-purpose web server serving data in the form of Prolog terms (possibly encoded in JSON) to clients as well as to other nodes. Just like a web server, a node has a domain name, makes use of HTTP and/or WebSockets, and is usually capable of serving more than one client at the same time. 

We use the term \textit{node} rather than \textit{server}, partly because this is what Erlang uses, but a more important reason is that a node is much \textit{more} than just a web server. A node is also a \textit{runtime system} for Web Prolog capable of executing Web Prolog programs stored by the node and/or passed to it by the client over the web APIs. The Prolog Web is a network of such code, ready to be executed by actors, ready to be queried by clients. 

\subsection{Web Prolog programs and data}

Data and programs stored at a node come in two flavors, characterized by two kinds of Web Prolog clauses. One flavor might have clauses looking like this:

\begin{verbatim}
mortal(Who) :- human(Who).
\end{verbatim}

\noindent Clauses of this flavor have logical interpretations, and this particular clause translates into the formula $\forall x [human(x) \rightarrow mortal(x)]$ in first-order predicate logic. 

As suggested in Figure~\ref{fig:prolog-web-intro}, a node can be queried from a terminal running in a web browser. Given our definition of \texttt{mortal/1} and a set of Prolog \textit{facts} such as

\begin{verbatim}
human(socrates).
human(plato).
human(aristotle).
\end{verbatim}

\noindent we can ask for the solutions to a query \texttt{?-mortal(Who)}. When we run the query from a terminal we enter into a kind of interaction with a Prolog process that most readers are probably familiar with -- in which the system only offers one solution at a time and the user may request further solutions by typing a semicolon:

\begin{verbatim}
?- mortal(Who).
Who = socrates ;
Who = plato ;
Who = aristotle.
?-
\end{verbatim}

\noindent Such queries can also be given logical interpretations, and this particular one translates into the formula $\exists x [mortal(x)]$. The interaction not only shows us that there exists an \textit{x} such as \textit{x} is mortal, but also provides us with examples. The point here is to remind ourselves that Prolog is a \textit{logic programming language} eminently suitable for \textit{knowledge representation}, or for other uses that require logic. The interaction also demonstrates that Prolog is a \textit{nondeterministic} programming language, meaning that a query may have several solutions, which can be lazily iterated over using Prolog's backtracking mechanism. 

In comparison with the clause defining the predicate \texttt{mortal/1}, the following clause has a very different flavor and looks a lot more like a program, and indeed something from the Erlang world -- a clause in the \textit{style} of Erlang:

\begin{verbatim}
echo_actor :-                     
    receive({                     
        echo(From, Msg) ->       
            From ! echo(Msg),   
            echo_actor  
    }).        
\end{verbatim}

\noindent As its name suggests, it implements the core of an echo server: an actor capable of  receiving messages and returning them to the sender as an ``echo.'' In contrast to the  definition of \texttt{mortal/1}, this clause cannot be given a logical interpretation; both \texttt{receive/1} and \texttt{!/2} are built-in predicates with side effects, so the clause must be  understood as a procedural construct rather than a declarative one -- a procedure that, when executed by an actor process, determines that actor's \textit{behavior}.

The predicate \texttt{spawn/2} creates a concurrent process running in parallel with the caller, binding the variable in its second argument to a process identifier:

\begin{verbatim}
?- spawn(echo_actor, Pid).
Pid = 10783471.
?-
\end{verbatim}

\noindent An interaction exercising the echo actor might look like this:

\begin{verbatim}
?- self(Self).
Self = 81902375.
?- 10783471 ! echo(81902375, hello).
true.
?- receive({Message -> true}).
Message = echo(hello).
?- 
\end{verbatim}

\noindent Readers familiar with Erlang will recognise the pattern. At this point other readers only needs to know that the predicate \texttt{self/1} retrieves the pid of the calling process, \texttt{!/2} sends a message to another process, and \texttt{receive/1} waits for an incoming message. Even this tiny example already conveys the essential properties of an actor. It has its own mailbox, reacts asynchronously to messages, can run indefinitely, and maintains identity through its process identifier. 

The spawn operation can be configured by means of \textit{options}, so there is a \texttt{spawn/3} as well, where the optional last argument is a list of options. For example, the \texttt{node} option takes a URI and allows an actor to be spawned on the remote node to which that URI points:

\begin{verbatim}
?- spawn(echo_server, Pid, [
       node('http://n4.org')
   ]).
Pid = 41983256@'http://n4.org'.
?- 
\end{verbatim}

\noindent This demonstrates that our actor model is not just for local concurrency but is inherently network-transparent and distributed -- this is the ``Web'' in ``Web Prolog.''

Other options inject source code into the spawned actor's private database, or specify what should happen when the actor dies, i.e. what its parent (if it has one) might learn about the reason for its death, and how its children (if it has any) will be treated when it terminates. These aspects of actor lifecycle management will be explored in greater depth later in the book.

Although the echo actor is intentionally simple -- a specialized process whose only task is to reflect incoming messages back to the sender -- it represents only the most basic form of actor behaviour. Through the combined use of options and meta-programming techniques, Web Prolog, like Erlang, supports a wide range of \textit{generic} behavioural patterns. These include actors that implement a generic client-server model, actors capable of actors supervising other actors, and actors that function as event-driven, hierarchical state machines. We shall return to these behavioural patterns when we discuss server templates, supervision trees, and statechart actors in later chapters -- including several powerful patterns that cannot easily be expressed in Erlang but arise naturally from properties that Prolog, but not Erlang, has.

\section{The current state of Prolog}

Created in 1972, Prolog is a rather old programming language. Any rumors of its imminent death can easily be debunked, however, as it is currently experiencing something of a revival. New systems are being implemented; features such as tabling and constraint-logic programming libraries are being added to both new and established systems; new books are being written; the community is small but thriving, and the language is still taught at some universities.

Still, the current state of Prolog is far from ideal. The community is \textit{too} small and lacks the cultural visibility that defines a healthy, growing language ecosystem. This is well known and was fairly recently documented in a series of papers.

In 2022, as part of the celebration of the Year of Prolog,\footnote{\url{https://prologyear.logicprogramming.org}} a number of logic-programming scholars published a survey paper titled \textit{Fifty Years of Prolog and Beyond} (hereafter referred to as FYPB), as well as a book of papers titled \textit{Prolog: The Next 50 Years} (henceforth PN50Y), which followed a year later. As they are written by a substantial portion of the \textit{cr\`eme de la cr\`eme} of the Prolog community, it makes sense to take these papers as probably the best and most current descriptions of the history of Prolog, its current state, and what the authors wish for its future. While this book does not discuss the history of Prolog, its \textit{current} state and its future are certainly relevant.

What makes FYPB and PN50Y suitable as background for the project described in this book -- and the reason we shall refer to them at various points going forward -- is that we share a common interest with their authors: to improve Prolog and its ecosystem as a whole. We ask ourselves the same questions they pose in the following passage from FYPB:

\begin{quote}
How might Prolog and its community evolve in the future? Can we better unify the new aspects that are offered by different implementations? How should efforts for increased portability be organized? Does it make sense to aim for a unified language? And what tools could be provided to ease development? Furthermore, we propose a plan for future steps that need to be taken to evolve Prolog as a language.
\end{quote}

\noindent However, as will become apparent, we arrive at different answers and at different ideas for how things might be improved. The Prolog Trinity ecosystem should be seen as our overarching proposal for these various ideas -- a scaffolding that helps ensure that they are mutually coherent and consistent, improves the cultural visibility of the current Prolog ecosystem, steadies its state, and helps heal what is broken.

\section{Motivating the Prolog Trinity ecosystem}

A vision may be coherent, a design elegant, and an implementation technically feasible, yet this alone does not suffice to justify its realization. What is also required is a clear motivation and a compelling argument for why the undertaking would be worthwhile.

The purpose of this book is to provide such an argument, and to show that its realization would not only benefit the implementers and users of existing Prolog systems, but also strengthen the Prolog ecosystem and the community as a whole. The following subsections outline the main motivations for developing the Prolog Trinity ecosystem.


\subsection{Why Web Prolog?}\label{sec:why-web-prolog}

\begin{enumerate}

\item \textbf{A simple model for concurrency-oriented programming.}
Web Prolog brings a clear and comprehensible model of concurrent programming to the Prolog world. Concurrency has long been a source of complexity within the Prolog community, and the ISO threads model offers only a fairly low-level solution. By contrast, extending Prolog with Erlang-style constructs yields an elegant and expressive multi-paradigm language -- one that can, in many respects, be more capable than traditional Prolog or Erlang alone. The actor model provides a high-level alternative for structuring concurrent programs, with semantics that are simple enough to teach and reason about. Another practical advantage is that Erlang-style concurrent programming in Web Prolog can be learned from the substantial literature produced by the Erlang community.

\item \textbf{An elegant model for distributed programming.}
Erlang introduced a remarkably clear and effective model for distributed programming, based on the principle that communication between processes should be independent of their physical location. Web Prolog adopts this foundation while embedding it in a logic programming context. It treats distribution as a natural extension of local message passing: actors communicate across nodes with the same simplicity and uniform semantics as within a single node. By preserving Erlang's conceptual clarity and pragmatic approach to fault handling, Web Prolog offers a transparent model of distributed computation -- one that combines the expressive reasoning of Prolog with the robust, process-oriented discipline of Erlang.

%\item \textbf{A profile of Prolog for web programming.}
%Although Prolog extended with Erlang-style concurrency primitives already forms a capable general-purpose language, Web Prolog focuses on a more specific goal: to define an accessible and coherent \emph{profile} of Prolog for programming the Web. Drawing on Erlang's distributed model yet reimagining it for an open, heterogeneous, and loosely federated environment, Web Prolog provides a foundation for scripting the Web with logic and implementing communication protocols among distributed agents. In other words, Web Prolog is intentionally bilingual at the transport level while remaining monolingual at the language level: the same Prolog terms and the same programming idioms scale from request--response querying to long-lived actor conversations. Two complementary modes of interaction are supported: declarative query answering over stateless or semi-stateful HTTP, and stateful actor communication over persistent WebSocket connections. Together these mechanisms cover both short-lived and enduring forms of interaction, making Web Prolog a versatile and elegant language for the programmable Web.

\item \textbf{A profile of Prolog for web programming.}
Although Prolog extended with Erlang-style concurrency primitives already forms a capable general-purpose language, Web Prolog focuses on a more specific goal: to define an accessible and coherent \emph{profile} of Prolog for programming the Web. Two complementary modes of interaction are supported: declarative query answering over stateless or semi-stateful HTTP, and stateful actor communication over persistent WebSocket connections. Drawing on Erlang's distributed model yet reimagining it for an open, heterogeneous, and loosely federated environment, Web Prolog provides a foundation for scripting the Web with logic and for implementing communication protocols among distributed agents. In other words, Web Prolog is intentionally bilingual at the transport level while remaining monolingual at the language level: the same Prolog terms and the same programming idioms scale from request--response querying to long-lived actor conversations. Together these mechanisms cover both short-lived and enduring forms of interaction, making Web Prolog a versatile and elegant language for the programmable Web.

\end{enumerate}

\subsection{Why Prolog agents?}\label{sec:why-prolog-agents}

\begin{enumerate}

\item \textbf{Agents as powerful programming abstractions.}
Prolog agents provide a simple yet expressive abstraction for continuous and interactive behaviour. A Prolog agent is a software process capable of communicating with other agents on the Prolog Web, exchanging Prolog terms of arbitrary complexity and, when appropriate, fragments of Web Prolog programs themselves. In this book we treat both actors and nodes as agents, and both come in specialised varieties such as toplevel actors, shells, statechart actors, and nodes with different profiles and capabilities. Because agents can be queried, instructed, supervised, and composed into larger systems, Web Prolog becomes not only a language for logic programming but also a language for \emph{agent programming}, supporting the development of single agents as well as broader multi-agent systems. The result is a unified abstraction for autonomous, interactive components at the centre of the ecosystem.

\item \textbf{Statechart actors.}
Statecharts, introduced by David Harel, offer a visual formalism for describing complex, reactive state machines with hierarchy, deterministic parallelism, and well-defined event semantics. Statechart actors are Prolog agents that provide a principled and highly expressive model for event-driven control. When viewed as actors, statechart processes become natural participants in a message-passing environment. Using Web Prolog as its data modelling and scripting language allows statechart actors to integrate declarative reasoning with robust, event-driven control. This combination yields a well-suited platform for constructing AI applications of many kinds, particularly intelligent conversational or embodied agents. Incorporating statechart actors broadens the ecosystem's ability to support sophisticated reactive systems. The point is not to add another agent type, but to make \emph{control} a first-class, inspectable artefact in the same ecosystem as declarative models and message-passing interaction.

\item \textbf{Turn-key Prolog nodes.}
An important practical motivation for the Prolog agent model is the availability of \emph{turn-key Prolog nodes} -- ready-made, nearly complete applications that become operational as soon as domain logic and configuration are supplied. Turn-key nodes lower the threshold for participation by capturing proven patterns for supervision, communication, persistence, and deployment, thereby sparing developers from boilerplate concerns and allowing them to focus on domain-specific reasoning and interaction. Because these nodes act as agents in their own right, they can be queried, extended, and orchestrated alongside user-defined agents, making the environment uniform and easy to understand. Such simplicity and accessibility are essential if the Prolog Web is to grow into a living shared environment rather than remain a conceptual construct.

\end{enumerate}

\subsection{Why the Prolog Web?}\label{sec:why-prolog-web}

\begin{enumerate}

\item \textbf{A network of logic-driven Prolog nodes.}
A central motivation for the Prolog Web is that a Prolog node is not a monolithic service but a participant in a distributed network of message-passing, logic-driven agents. Turn-key Prolog nodes provide the structure needed to populate this network: self-contained, declarative services that embody the communication, reasoning, and supervision patterns of Web Prolog. Nodes share a uniform design so that newcomers can learn them quickly, practitioners can use them effectively, and deployers can operate them with minimal effort. This consistency allows the Prolog Web to grow organically into a coherent, agent-populated environment rather than a collection of isolated services.

\item \textbf{A natural foundation for modern AI.}
Viewed in this light, the Prolog Web is not merely an attractive programming environment but a natural foundation for building modern AI systems. Logic-based, message-passing agents are well suited to the construction of hybrid and classical multi-agent systems, distributed reasoning frameworks, and embodied conversational agents that require explainability, coordination, and long-lived interaction. The Prolog Trinity ecosystem aligns these requirements with web-native infrastructure, providing a principled model for constructing AI that is transparent, distributed, and inherently capable of reasoning. This makes the Prolog Web an appealing platform for developing the emerging \emph{Agentic Web} (a buzzword for a web of autonomous, reasoning agents).

\item \textbf{Nondeterministic RPC as a web-programming capability.}
Another motivation for the Prolog Web is that it supports forms of distributed interaction that are difficult to realise cleanly in other web programming models. In particular, its notion of nondeterministic remote procedure calls allows a client to request multiple answers to a query from a remote node, mirroring Prolog's local search behaviour across the network. This preserves the declarative reading of logic programs while extending it seamlessly to distributed settings, enabling agents to cooperate in solving search problems, generating alternative plans, or exploring hypotheses in parallel. Treating nondeterministic RPC as a first-class web capability makes the Prolog Web a distinctive and expressive environment for distributed logic programming.

\item \textbf{Federation with explicit boundaries.} 
The Prolog Web is designed as a loosely coupled federation: nodes are free to specialise internally, but must expose small, well-specified interfaces at the boundary. This matches how the Web scales, and it also creates a natural place for capability restrictions, resource limits, and ownership policies -- concerns that become unavoidable once we take openness seriously.

\end{enumerate}

\subsection{Why the Prolog Trinity ecosystem as a whole?}\label{sec:why-trinity-ecosystem}

Figure~\ref{fig:prolog-web-intro} does not reveal anything about the \textit{implementation} of the nodes it depicts. They could all be implemented on top of a single Prolog system, or several systems may be involved, perhaps one for each node, reflecting the diversity of the Prolog landscape. Building the Prolog Web on top of one system is probably the fastest path to a public demo, early feedback, and something that runs on today's Web and demonstrates value. In a research-oriented community such as Prolog's, stopping at that stage would not be unusual. Our proposal, however, is to turn the endeavour into a collective, community-driven project.

The Prolog Trinity ecosystem is therefore motivated by the need for a common scaffolding that supports, rather than suppresses, the existing diversity of Prolog systems. Different implementations excel in different domains: some prioritise performance, others constraint solving, tabling, or integration with foreign languages. Rather than seeking unification through a large monolithic standard or a single system, we treat this diversity as a resource. Web Prolog is defined as a small, implementable, web-aligned fragment that acts as a \emph{lingua franca} for cross-system interoperability. By wrapping each system in a node that speaks Web Prolog, heterogeneous engines can collaborate in real time without surrendering their strengths. This is not a call for unification by portability alone; it is a call for precise, web-aligned interfaces among nodes using small, well-specified protocols. Modest systems can implement just the essential subset and still participate fully. The result is a federation that mirrors how the Web itself grows: loosely coupled, strictly specified at the boundaries, and free to specialise behind them.

To move from the narrow view of a single implementation to the broader perspective of the Prolog ecosystem, we must also relate our aims to concerns already articulated by the community. As we learned above, FYPB and PN50Y ask how Prolog might evolve, how portability might be improved, whether a unified language is sensible, and what tools would best support development. These questions form a shared point of departure, yet our proposed answers differ in important ways. Where FYPB sees weakness in ``limited portability'', we concentrate on a portable, carefully defined \emph{fragment} of Prolog as a basis for interoperability. Where they see a problem of \emph{fragmentation}, we see an opportunity for \emph{diversification}, tied together by a clear Web Prolog substrate. Where they identify a threat in the ``post-desktop world of JavaScript web applications'', we see an opportunity to design a logic-based, actor-oriented, federated programming environment capable of thriving on the Web itself.

Our programme complements the forward-looking agendas in FYPB and PN50Y: where they outline desirable directions for the future of Prolog, we propose a web-native and agent-centric realisation path, complete with concrete artefacts that can be implemented, tested, compared, and incrementally standardised. For these reasons, we argue that the realisation of the Prolog Trinity ecosystem is not merely feasible but worthwhile: it offers simplicity for developers, coherence for the community, and a principled platform for the next generation of AI on the Web.

Finally, a note on scope. The Trinity is not proposed as a replacement for existing Prolog systems, but as a way to make them cooperate. It is not a demand for a monolithic unified Prolog, but an attempt to define small boundaries that many systems can implement. And it is not a promise of a frictionless open network; rather, it is an argument that openness becomes manageable once capabilities and interfaces are explicit.


\section{Webizing Prolog and prologizing the Web}

\begin{quotation}
\noindent The essential process in webizing is to take a system which is
designed as a closed world, and then ask what happens when it is
considered as part of an open world. Practically, this effect on a
computer language is to replace the names/tokens/identifiers for URIs.
Thus, where before reference could only be made to something in the
same document/program/module one can with equal ease make reference to
something in a different one somewhere in that abstract space which is
the Web.
\vspace{-0.1cm}
\flushright  \textit{Tim Berners-Lee, 1998}
\vspace{0.1cm}
\end{quotation}

\noindent Some thirty years in the making, the World Wide Web is a formidable success and the biggest distributed programming system ever constructed, and has, over the years, become a key delivery platform for a variety of sophisticated interactive applications in just about any conceivable domain. 

The traditional web, already in its most primitive form as a Web of Documents (also known as Web 1.0), is \textit{distributed} (running on more than one machine), \textit{decentralized} (running at different geographical locations and/or having different owners) and \textit{open} (anyone can use and contribute). These are some of the traits the design of Prolog Web seeks to inherit from the conventional web. Whereas distribution comes with the territory, decentralization and openness have more to do with \textit{web culture}. 

\subsection{Webizing Prolog}

To develop the design of the Prolog Trinity ecosystem we need to apply the process of \textit{webizing} to Prolog. Here is what we need to do, and each of the points must be meticulously addressed:

\begin{enumerate}
\def\labelenumi{\arabic{enumi}.}
\item
  \textbf{Introduce URIs into the ecosystem}: Uniform Resource
  Identifiers (URIs) are fundamental to the web as they uniquely
  identify resources. Integrating URIs into Prolog will enable the
  language to directly address and interact with the plethora of
  resources available on the web. 
\item
  \textbf{Exploit the existing web infrastructure}: The Web's infrastructure includes protocols such as HTTP and WebSocket as well as data formats such as JSON. Web Prolog should be able to utilize these protocols and data formats seamlessly, and as much as possible behind the scenes.
\item
  \textbf{Make use of existing means for security}: Web security is a
  multifaceted domain, encompassing aspects like authentication,
  authorization, data integrity, and confidentiality. Web Prolog needs
  to integrate with existing security mechanisms such as TLS for
  encrypted communications, OAuth for authorization, and possibly more
  sophisticated web-specific security protocols. These mechanisms must
  be ingrained within the Web Prolog's runtime environment to protect
  against common web vulnerabilities.
\item
  \textbf{Ensure web-size scalability}: The scalability challenge for
  Web Prolog is twofold: it must handle large numbers of simultaneous
  connections and manage extensive datasets. This will require efficient multitasking and concurrency
  models, optimized algorithms for distributed computing, and an
  architecture that allows for load balancing and a lot of caching.

\item
  \textbf{Support web application development}: To truly support web
  application development, Web Prolog should offer a
  development experience that is both powerful for experienced
  developers and accessible to those new to Prolog.
\item
  \textbf{Ensure fitness to web culture}: The web has a distinct culture
  that prioritizes openness, collaboration, and community-driven
  development. Web Prolog should embrace this by being open source,
  having a clear and liberal licensing model, providing excellent
  documentation and support, and fostering a healthy community around it. 
\item
  \textbf{Aim for standardization}: Finally, to gain widespread
  adoption, Web Prolog must strive for standardization. This entails
  working with standard bodies and the broader community to establish an
  official specification for Web Prolog. Standardization helps ensure
  interoperability between different implementations and provides a
  stable foundation upon which developers can build.
\item
  \textbf{Play nicely with existing web standards}: For Prolog to be web-friendly, it must enter into fruitful relations with existing W3C standards. This implies that Web Prolog must be able to handle standard web protocols and document formats, work efficiently with web APIs, and follow guidelines such as for internationalization. Prolog\textquotesingle s strength in dealing with ontologies and rule-based systems allows it to play nicely with the existing standards for the Semantic Web. We believe we may also be able to connect to a W3C's SCXML, although this may be harder. More on this later.
\end{enumerate}

\subsection{Prologizing the Web}

\noindent Webizing Prolog involves adapting and extending Prolog to operate effectively in the context of the Web. The focus is on Prolog and the task of designing and developing a special-purpose profile that makes it more useful for web programming than it currently is. The phrase ``prologizing the Web'' is just another way to describe the very same process, aiming at bringing Prolog and the Web closer together. Prologizing the Web involves extending the traditional web with the capabilities of Prolog. We put the Web in focus and seek to improve and further evolve it by the addition of Prolog. We do this because we believe that a prologized Web is better than a Web not prologized.  

\begin{enumerate}
\def\labelenumi{\arabic{enumi}.}
\item
  \textbf{Ensure clever use of Prolog's capabilities}: Make sure that the capabilities that Prolog has, such as knowledge representation, reasoning, problem solving and parsing can be put to good use on the Prolog Web. 
\item  
  \textbf{Leverage the power of Erlang-style concurrency\textbf{}:} Ensure that the actor-based
  predicates for process creation and message passing can be used to support
  concurrency and distribution across the Web.
\item
  \textbf{Ensure that Web Prolog runs in web browsers}: Do not introduce constructs that make it impossible to implement Web Prolog in browsers.
\end{enumerate}

\noindent Towards the end of the book, in Chapter~\ref{sec:tech-assesment}, we shall return to these points, and assess what has been done to webize Prolog, prologize the Web and create a design for the infrastructure supporting the Prolog Trinity ecosystem.

\section{Aiming for a standard}

The exclamation mark at the end of the book's subtitle might lead some readers to suspect that it can be read as something of a manifesto or mission statement. There is indeed some truth in that, as we are trying to draw the intellectual support necessary for an attempt to create \textit{web standards} for the technologies relevant to the realization of the Prolog Trinity ecosystem. Presumably, such an attempt can only succeed if supported by a reasonably large number of individuals from the relevant communities -- by people from the Prolog and Erlang communities, people in the Semantic Web community, and people interested in AI and in building conversational systems.


\section{Thirty years of Prolog on the Web}\label{sec:prolog-on-the-web}

We do not want to give the impression that the suitability of Prolog for programming the Web has not been investigated before. On the contrary, in the mid-1990s the advent of the Web sparked a great deal of interest in the logic programming community, reflected in two workshops devoted to the subject and in the development of libraries that mostly relied on creating so-called CGI programs in Prolog. Here is a snippet of text lifted from the call for papers for the second workshop:\footnote{\url{http://www.cliplab.org/lpnet/proceedings97/preface.html}}

\begin{quote}
This workshop is the second in a series intended to explore the elective affinities between Logic Programming and Internet technologies with emphasis on enhancing the World Wide Web with knowledge, deductive abilities and superior forms of interactive behavior. With the paradigm shift to highly inter-connected computers and programming tools, logic programming languages have a unique opportunity to contribute to practical Internet application development. Simplicity, remote executability, robustness, automatic memory management, are among the features some LP languages share with emerging tools like Java. Superior meta-programming and high-level distributed programming facilities, built-in grammars and dynamic databases, declarative semantics are among their competitive advantages.
\end{quote}

\noindent Since then the Web has developed at an amazingly fast pace, and web applications have become increasingly speedier and more interactive. The Web of today is \textit{very} different from what it was almost thirty years ago -- faster, more reliable, more mobile, and with a fantastic selection of browser APIs that lets a web developer add features such as spoken interaction, video or 3D graphics to an application. There is a reason we nowadays talk about, not only the Semantic Web, but also the Conversational Web, the Mobile Web, the Web of Things, and (more recently) the Agentic Web.\footnote{See \cite{lagermyrendal2012} for a description of the many facets of the current Web.} Furthermore, many (potentially) disruptive technologies such as automation, artificial intelligence, blockchain technology, wearable technology, augmented reality, virtual reality, 5G, 3D printing, robotics and quantum computing are likely to more or less depend on and be integrated into the Web. In this book we focus on artificial intelligence and its integration into the Web.

With the Web (technically) so much better and so different from what it used to be, there is, in our opinion, every reason to make a new attempt at ``enhancing the World Wide Web with knowledge, deductive abilities and superior forms of interactive behavior.'' In particular, ``superior forms of interactive behavior'' is likely to depend on the WebSocket protocol and/or the implementation of Web Prolog in browsers or (mobile) apps. The WebSocket standard did not exist twenty years ago, which meant that every approach to web logic programming had to work solely with HTTP. While HTTP remains the most important protocol (and as we shall see, two out of the three APIs driving the Prolog Trinity ecosystem rely on it), important aspects of some of the inventions described in this book rely on the WebSocket protocol, a so called server-push technology which opens a whole slew of opportunities for developers of web applications based on Prolog. 

Regarding the history of Prolog, we do not suggest that nothing has happened between now and then.
From the 1990s onward, Prolog has repeatedly adapted to the web's changing landscape. Early work such as the CLIP group's PiLLoW library showed how HTML and HTTP could be handled declaratively by treating web structures as Prolog terms. With the rise of the Semantic Web in the 2000s, Prolog gained renewed relevance: SWI-Prolog's RDF store, web libraries, and the ClioPatria platform demonstrated that a Prolog system could serve simultaneously as a high-performance knowledge base and a full HTTP server.

In the 2010s, more interactive and distributed styles of web programming led to further innovations. Pengines (Lager \& Wielemaker) made server-side Prolog engines accessible from JavaScript through a simple RPC model, enabling incremental solution streaming and smooth integration with modern web applications. At the same time, projects such as Tau Prolog brought ISO-style Prolog directly into the browser, allowing entirely client-side logic applications written in JavaScript.

More recently, WebAssembly has made it possible to run full Prolog systems inside the browser. SWI-Prolog can now be compiled to WASM via Emscripten, and lightweight systems such as Scryer and Trealla have been ported as well. These provide fast, sandboxed Prolog engines that interoperate smoothly with JavaScript, opening the way for high-performance logic programming in modern web applications.

Among Prolog systems, SWI-Prolog in particular has all the right tools for building server-side support for web applications, in the form of mature libraries for protocols such as HTTP and WebSocket, and excellent support for formats such as HTML, XML and JSON. Therefore, it might be argued that since Prolog can and has been used for building server-side support for web applications, it should already be counted as a web programming language. But since this is true also for Python, or Ruby, or Java, or just about any other general-purpose programming language in existence, it would make the notion of a web programming language more or less empty.

We could of course argue that SWI-Prolog is much better at providing what is needed, but we shall go further than that. We believe that Prolog should claim an even more prominent position in this space. Web Prolog should be used not only for building server-side support for web applications, but should be considered as \textit{part} of the Web, in much the same way as HTML, CSS, JavaScript, RDF and OWL are parts of the Web. In other words, as we think it deserves it, we would like it to be recognized as \textit{a web technology in its own right}.


\section{The structure of the book, and a pointer to its home page}

In the hope that it might help the reader navigate the book's most important ideas and their connections, the mindmap in Figure~\ref{fig:summary-diagram-first} expands the exposition of the three most important aspects of the Prolog Trinity ecosystem.

\begin{figure}[h]
    \centering
	\includegraphics[width=\textwidth]{summary-diagram-first}
    \vspace{0.1cm}\caption{Mind map of the Prolog Trinity ecosystem and its main conceptual threads.}
    \label{fig:summary-diagram-first}
\end{figure}

\subsection{The structure of the book}

The mindmap reflects the structure of the rest of this book as follows. 


\subsection*{Part I: The conceptual and architectural core}

\textbf{Chapter~\ref{sec:web-prolog}} introduces Web Prolog as a Prolog-based, multi-paradigm language for concurrent and distributed programming. It explains what is inherited from standard Prolog, what is added to support actor-style concurrency, and why Erlang is a particularly useful reference point. To make the comparison concrete, the chapter translates a range of Erlang programs into Web Prolog and uses these translations to clarify the intended operational behaviour.

\medskip
\noindent \textbf{Chapter~\ref{sec:prolog-agents}} develops the agent strand of the mindmap. It introduces a simple taxonomy in which Prolog actors and Prolog nodes are the primary categories, and it focuses on the two agent forms that matter most for the rest of the book: the Prolog toplevel actor and the Prolog node. The chapter describes the built-in interaction protocol that distinguishes a toplevel actor from other actors, breaks a node down into its main components (controller, Web APIs, and the node-resident Prolog database), and illustrates the intended client--node interaction patterns through a series of scenarios. Statechart actors are treated separately in Chapter~\ref{sec:web-prolog-statechart-behavior}.

\medskip
\noindent \textbf{Chapter~\ref{sec:the-prolog-web}} presents the architectural proposal for the Prolog Web and shows how multiple nodes can be composed into an open, decentralised network. A central theme is distributed programming with an emphasis on programming models that aim to be close to network-transparent when appropriate. The chapter introduces \texttt{rpc/2-3} as a user-facing construct for nondeterministic remote procedure calls and uses it to delineate a sequential partition of the Prolog Web alongside the concurrent, mailbox-driven partition. It also makes explicit several axes along which designs can vary, including concurrency vs.\ sequentiality and purer vs.\ less pure interaction styles.

\subsection*{Part II: Behaviors and other generics}

\textbf{Chapter~\ref{sec:implementing-prolog-generics}} implements a set of generic constructs that capture reusable actor patterns. These generics can be specialised or configured after creation-time and encapsulate recurring concerns such as message protocols, supervision, and fault handling.

\medskip
\noindent \textbf{Chapter~\ref{sec:web-prolog-statechart-behavior}} introduces an additional category of actor, termed the \textit{statechart actor}, which represents an advanced form of event-driven state machine. The distinctive feature of a statechart actor, as compared to a conventional actor, is its implementation in a variant of an XML-based statechart language. In this paradigm, Web Prolog is employed solely for data modeling and scripting purposes, rather than serving as the primary programming language for the statechart logic.

\subsection*{Part III: The broader technological contexts}

\textbf{Chapter~\ref{sec:wp-ws-sw}} presents a detailed comparative analysis between the Prolog Web and the Semantic Web. In this chapter, we contend that the Prolog Web can be effectively constructed on top of the Semantic Web. Both frameworks are fundamentally grounded in computational logic; however, the Prolog Web, viewed as a relational logic-based environment, exhibits a natural affinity for integration with the Semantic Web. This intrinsic compatibility positions the Prolog Web as a promising candidate for Semantic Web programming, thereby enabling the development of applications that leverage the strengths of both paradigms.

\medskip
\noindent \textbf{Chapter~\ref{sec:prolog-ai-agents-on-the-agentic-web}} turns to AI applications, with particular attention to conversational and potentially embodied agents operating on the Web. It uses this setting to explore how statechart-style control structure and Web Prolog-style symbolic scripting can be combined in a pragmatic architecture for building interactive, tool-using agents.

\subsection*{Part IV: What's in it for you and your community?}

\textbf{Chapter~\ref{sec:whats-in-it-for-the-prolog-community}} asks what the Trinity offers to individual users and to the Prolog community as a whole. It frames the contribution in terms of practical capabilities, interoperability, and ecosystem-level opportunities, and it discusses how some commonly identified risks for Prolog might be mitigated.

\medskip
\noindent \textbf{Chapter~\ref{sec:for-other-communities}} asks the corresponding question for neighbouring communities, including those approaching the problem from distributed systems, Web engineering, and AI perspectives.

\subsection*{Part V: Final words}

\textbf{Chapter~\ref{sec:from-vision-to-execution}} moves from vision to execution by outlining an implementation path, identifying dependencies and engineering constraints, and mapping the proposal onto a concrete roadmap.

\subsection{The book's home page}

The material presented here is accompanied by a home page for the Prolog Trinity ecosystem, where updated code, examples, draft specifications, and links to related projects will be made available. A link to the home page is here (but not yet...):

\begin{verbatim}
http://trinity.elfenbenstornet.se
\end{verbatim}


%\section{The structure of the book, and a pointer to its home page}
%
%In the hope that it might help the reader navigate the book's most important ideas and their connections, the mindmap in Figure~\ref{fig:summary-diagram-first} expands the exposition of the three most important aspects of the Prolog Trinity ecosystem.
%
%\begin{figure}[h]
%    \centering
%	\includegraphics[width=\textwidth]{summary-diagram-first}
%    \vspace{0.1cm}\caption{A mindmap over the Prolog Trinity ecosystem. }
%    \label{fig:summary-diagram-first}
%\end{figure}
%
%\noindent The mindmap reflects the structure of the rest of this book as follows. Chapter~\ref{sec:web-prolog} presents the language of Web Prolog, its roots in standard Prolog, and its relation to the Erlang programming language. We explain why we regard Web Prolog as a multi-paradigm language, and why we believe that it has great potential for serving as a scripting language for programming the Web. With the goal of seeing how closely it can align with Erlang-style concurrent programming while remaining fundamentally Prolog, we undertake the translation of a fair number of Erlang programs into Web Prolog and demonstrate that they execute just as they would do in Erlang.
%
%Chapter~\ref{sec:prolog-agents} deals with agents, and with Prolog agents in particular. As suggested by the diagram, they come in different shapes and sizes, and can be arranged into a simple taxonomy where Prolog actors and Prolog nodes represent the primary subcategories. While Chapter~\ref{sec:web-prolog} will give many concrete examples of actors (represented by the ... in Figure~\ref{fig:summary-diagram-first}), and Chapter~\ref{sec:prolog-agents} will briefly revisit the abstract concept of the Prolog actor, most attention is devoted to Prolog toplevel actors and nodes. (Statechart actors are covered in a separate chapter.) In particular, the built-in communication protocol that distinguishes the Prolog toplevel from other actors is described. The Prolog node is broken down into its key components, such as the node controller, the Web APIs, and the node-resident Prolog database holding programs and data shared by actors running on the node. Specific scenarios will be used to demonstrate how the interaction between user clients and nodes is intended to work.
%
%Chapter~\ref{sec:the-prolog-web} presents our architectural proposal for the Prolog Web. A central theme is \textit{distributed programming}, with particular emphasis on (almost) \textit{network-transparent} distributed programming -- another strength of Erlang that has inspired us. In this chapter, we provide code examples that demonstrate inter-node communication within the Prolog Web and explore several scenarios in detail to illustrate its intended operation. Additionally, we introduce a user-friendly built-in predicate, \texttt{rpc/2-3}, which facilitates Prolog-style \textit{nondeterministic remote procedure calls}. We also delineate the various partitions of the Prolog Web, explicitly distinguishing between concurrent and sequential partitions, as well as between pure and less pure partitions. As illustrated in the diagram, we consider the Prolog Web to be an extension of the traditional Web and we aim for it to remain as decentralized, open, and scalable as the traditional Web itself.
%
%
%Chapter~\ref{sec:implementing-prolog-generics} implements a set of generic constructs that capture reusable actor patterns. These generics can be specialised or configured after creation-time and encapsulate recurring concerns such as message protocols, supervision, fault handling, and state management.
%
%Chapter~\ref{sec:web-prolog-statechart-behavior} introduces an additional category of actor, termed the \textit{statechart actor}, which represents an advanced form of event-driven state machine. The distinctive feature of a statechart actor, as compared to a conventional actor, is its implementation in a variant of an XML-based language known as State Chart XML (SCXML). In this paradigm, Web Prolog is employed solely for data modeling and scripting purposes, rather than serving as the primary programming language for the statechart logic.
%
%Chapter~\ref{sec:wp-ws-sw} presents a detailed comparative analysis between the Prolog Web and the Semantic Web. In this chapter, we contend that the Prolog Web can be effectively constructed on top of the Semantic Web. Both frameworks are fundamentally grounded in computational logic; however, the Prolog Web, viewed as a relational logic-based environment, exhibits a natural affinity for integration with the Semantic Web. This intrinsic compatibility positions the Prolog Web as a promising candidate for Semantic Web programming, thereby enabling the development of applications that leverage the strengths of both paradigms.
%
%Chapter~\ref{sec:prolog-ai-agents-on-the-agentic-web} examines the domain of artificial intelligence with a particular focus on embodied conversational agents (ECAs). We posit that ECAs are poised to become the preeminent face of AI, with the Web emerging as the primary arena for their widespread adoption and seamless integration. In this chapter, we further explore the synergistic potential of combining SCXML with Web Prolog, arguing that this integration provides a robust platform for constructing a diverse array of AI applications. Specifically, we emphasize the suitability of this hybrid approach for developing intelligent -- and potentially embodied -- conversational agents, thereby illustrating a viable pathway for advancing modern AI technologies.
%
%In Chapter~\ref{sec:whats-in-it-for-the-prolog-community} we ask ourselves: What concrete benefits does the Trinity offer to individual users and to the wider Prolog community? Our aim is to articulate not just what the Prolog Trinity \emph{is}, but what it \emph{does} -- what problems it solves, what new horizons it opens, and how it can empower the Prolog community to build a more intelligent, interconnected, and logical Web. We explain how some of the threats observed by the authors of FYPB might be mitigated, and we explore some of the opportunities provided by the Prolog Trinity ecosystem and the interoperability between agents.
%
%Chapter~\ref{sec:for-other-communities} asks the corresponding question for neighbouring communities, including those approaching the problem from distributed systems, Web engineering, and AI perspectives.
%
%Chapter~\ref{sec:from-vision-to-execution} moves from vision to execution by outlining an implementation path, identifying dependencies and engineering constraints, and mapping the proposal onto a concrete roadmap.
%
%===
%
%In Chapter~\ref{sec:steps-towards-standardisation} we will say more about standardization and the steps that need to be taken for it to succeed. We discuss the possibility of structuring the work as a project for which funding might be sought. We also argue that it might make sense to standardize Web Prolog and to develop a version 2.0 of SCXML.
%
%Chapter~\ref{sec:profiles-and-implementation} is about the implementation of a node. Snippets of code written in SWI-Prolog are provided that implement some of the central pieces of a Web Prolog system.
%
%
%
%Appendices A--E provide reference materials. Appendix~\ref{sec:manual} defines core APIs and predicates of Web Prolog, including concurrency, toplevels, and RPC.
%
%Appendix~\ref{sec:benchmarking} presents performance benchmarks, comparing Web Prolog to Erlang and analysing design trade-offs.
%
%Appendix~\ref{sec:big-example} provides a full concurrent example -- the Dining Philosophers problem -- demonstrating actor design and lifecycle coordination.
%
%Appendix~\ref{sec:web-api-examples} showcases client-side interaction with Web Prolog nodes using REST and WebSockets.
%
%Appendix~\ref{sec:teaching-material} lists pedagogical references for learning Prolog and Erlang, offering a curricular bridge into the system.
%
%The book moves from high-level ideas to concrete designs and code, and it ends by inviting others to use, adapt, and extend the work. Many chapters conclude with short \emph{discussion sections}. Their purpose is to step back from the technical material, reflect on design choices, relate the proposals to broader themes in logic programming and the philosophy of computing, and suggest questions or directions for readers who wish to think further or contribute to the ongoing work.


